% Figure 2: one Self-Harness optimization round as a loop, with pictogram panels.
% Four stages: weakness mining -> independent candidate proposal -> screening
% tournament -> full validation of the winner, plus the two bounded fallbacks
% (re-attribution after an all-rejected screen; one-step refinement after a
% rejected winner).
% Relies on the palette (editgreen, editfill, fixgrey, fixfill, envtint, subtint,
% inkcol), card macros (\cardhd, \cardsb), and the penbadge/lockbadge pics defined
% in the preamble of aaai2027_paper.tex.
\definecolor{rejred}{HTML}{B3363D}
\definecolor{minehd}{HTML}{2B5FAD}
\definecolor{valhd}{HTML}{B26A05}
\definecolor{scrhd}{HTML}{6B4FA8}
\definecolor{scrtint}{HTML}{F1EDF8}
%
% ---- glyph library (figure-local) -------------------------------------------
% Nested tikzpictures inherit nothing from the outer picture; glyphrow restores
% fonts, colors, and arrow tips for the in-box panels.
\tikzset{
  % sharp corners: the enclosing stage/harnessbox node's rounded corners=7pt
  % stays in force inside the node text's nested tikzpicture and would turn
  % every cornered path (checkmarks, cells) into a smooth swoosh.
  glyphrow/.style={font=\sffamily\scriptsize, line width=0.6pt, inkcol,
    sharp corners, >={Stealth[round,length=3pt]}, every node/.style={text=inkcol,
    inner sep=1pt, align=center}, baseline=(current bounding box.center)},
  % task card: a mining-data task instance (plain document)
  taskcard/.pic={
    \draw[rounded corners=1.5pt, fill=white, draw=black!40] (-0.42,-0.5) rectangle (0.42,0.5);
    \foreach \y in {0.3,0.1,-0.1,-0.3}{
      \draw[line width=0.5pt, black!40] (-0.26,\y) -- (0.28,\y);}
  },
  % Execution call tree; uniform nodes carry no semantic correctness labels.
  calltree/.pic={
    % edges first, dots after, so the dots sit on top
    \draw[black!55] (0,0.55) -- (-0.45,0.0);
    \draw[black!55] (0,0.55) -- (0,0.0);
    \draw[black!55] (0,0.55) -- (0.45,0.0);
    \draw[black!55] (0.45,0.0) -- (0.27,-0.48);
    \draw[black!55] (0.45,0.0) -- (0.63,-0.48);
    \fill[inkcol] (0,0.55) circle (0.07);
    \fill[inkcol] (-0.45,0.0) circle (0.07);
    \fill[inkcol] (0,0.0) circle (0.07);
    \fill[inkcol] (0.45,0.0) circle (0.07);
    \fill[inkcol] (0.27,-0.48) circle (0.07);
    \fill[inkcol] (0.63,-0.48) circle (0.07);
  },
  % two failure clusters grouped by signature
  clusterpair/.pic={
    \fill[rejred] (-0.38,0.1) circle (0.07);
    \fill[rejred] (-0.52,-0.12) circle (0.07);
    \fill[rejred] (-0.24,-0.14) circle (0.07);
    \draw[rejred, dashed, line width=0.5pt] (-0.38,-0.04) ellipse (0.3 and 0.3);
    \fill[valhd] (0.38,0.12) circle (0.07);
    \fill[valhd] (0.24,-0.12) circle (0.07);
    \fill[valhd] (0.52,-0.1) circle (0.07);
    \draw[valhd, dashed, line width=0.5pt] (0.38,-0.02) ellipse (0.3 and 0.3);
  },
  % the fixed model acting as proposer; the H_t subtitle (a node placed in the
  % panel, not here -- nodes inside pics can drift off the pic's paths) marks
  % the current harness it edits
  proposercard/.pic={
    \draw[rounded corners=2pt, fill=white, draw=black!40] (-0.45,-0.68) rectangle (0.45,0.55);
    \fill[inkcol] (0,0.22) circle (0.11);
    \fill[inkcol] (-0.2,-0.18) arc (180:0:0.2) -- cycle;
    \path (0.45,0.55) pic[scale=1.3] {lockbadge};
  },
  % three stacked task cards: the small screening subset V_s
  subsetcards/.pic={
    \foreach \d in {0.16,0.08,0}{
      \draw[rounded corners=1.5pt, fill=white, draw=black!40]
        (-0.32+\d,-0.42+\d) rectangle (0.32+\d,0.38+\d);}
    \foreach \y in {0.2,0.02,-0.16}{
      \draw[line width=0.5pt, black!40] (-0.18,\y) -- (0.2,\y);}
  },
  % Winner and incumbent on the same full held-out validation tasks.
  splitcheck/.pic={
    \draw[rounded corners=1.5pt, fill=white, draw=black!40] (-0.45,0.06) rectangle (0.45,0.62);
    \draw[line width=0.5pt, black!40] (-0.3,0.44) -- (0.18,0.44);
    \draw[line width=0.5pt, black!40] (-0.3,0.24) -- (0.28,0.24);
    \draw[editgreen, line width=0.8pt, line cap=round] (0.2,0.4) -- (0.27,0.33) -- (0.4,0.5);
    \draw[rounded corners=1.5pt, fill=white, draw=black!50] (-0.45,-0.62) rectangle (0.45,-0.06);
    \draw[line width=0.5pt, black!40] (-0.3,-0.24) -- (0.18,-0.24);
    \draw[line width=0.5pt, black!40] (-0.3,-0.44) -- (0.28,-0.44);
    \draw[editgreen, line width=0.8pt, line cap=round] (0.2,-0.28) -- (0.27,-0.35) -- (0.4,-0.18);
  },
  % promotion decision marks
  acceptreject/.pic={
    \fill[editgreen] (0,0.24) circle (0.16);
    \draw[white, line width=1.1pt, line cap=round] (-0.08,0.24) -- (-0.02,0.16) -- (0.1,0.34);
    \fill[rejred] (0,-0.24) circle (0.16);
    \draw[white, line width=1.1pt, line cap=round] (-0.08,-0.16) -- (0.08,-0.32);
    \draw[white, line width=1.1pt, line cap=round] (-0.08,-0.32) -- (0.08,-0.16);
  },
  % Surface-grid cells (S1-S9, reading order). Each cell is ONE node whose own
  % border is the cell frame and whose text is the label -- separate \draw
  % rectangles + \node labels drift apart in this document (path-vs-node offset
  % triggered by earlier nested pictures), and a single node cannot misalign
  % with itself.
  cellnode/.style={rounded corners=1pt, inner sep=0pt, minimum width=0.38cm,
    minimum height=0.38cm, font=\sffamily\scriptsize, align=center},
  cellplain/.style={cellnode, fill=white, draw=fixgrey, line width=0.5pt, text=black!50},
  celledited/.style={cellnode, fill=editgreen!22, draw=editgreen, line width=0.9pt, text=editgreen},
  cellnew/.style={cellnode, fill=editgreen, draw=editgreen, line width=0.9pt, text=white},
  % unlabeled mini cells for the per-candidate surface grids
  minicell/.style={rounded corners=0.6pt, inner sep=0pt, minimum width=0.17cm,
    minimum height=0.17cm, fill=white, draw=fixgrey, line width=0.4pt},
  minihit/.style={minicell, fill=editgreen, draw=editgreen},
  % screening-rank bars and per-length delta bars (single nodes, see above)
  rankbar/.style={inner sep=0pt, minimum height=0.13cm, anchor=west},
  dbar/.style={inner sep=0pt, minimum width=0.2cm},
}
% \minigrid{x}{y}{hit}{label}: a 3x3 + 1 unlabeled surface grid whose cell
% number <hit> (1-10) is the candidate's single edited surface.
\newcommand{\minigrid}[4]{%
  \foreach \j [count=\r from 0] in {2,1,0}{
    \foreach \i in {0,1,2}{
      \pgfmathtruncatemacro{\slab}{\r*3+\i+1}
      \ifnum\slab=#3 \tikzset{c/.style=minihit}\else\tikzset{c/.style=minicell}\fi
      \node[c] at ({#1+\i*0.21},{#2+\j*0.21}) {};}}
  \ifnum#3=10 \tikzset{c/.style=minihit}\else\tikzset{c/.style=minicell}\fi
  \node[c] at ({#1+0.21},{#2-0.21}) {};
  \node[font=\sffamily\scriptsize] at ({#1+0.21},{#2-0.48}) {#4};
}
%
% AAAI: the three-column loop diagram spans both columns (figure*, [t] only).
\begin{figure*}[t]
\centering
\resizebox{\textwidth}{!}{%
\begin{tikzpicture}[
  font=\sffamily\small, line width=0.7pt, inkcol,
  >={Stealth[round,length=5pt]},
  every node/.style={text=inkcol},
  stage/.style={rounded corners=7pt, draw=black!30, line width=1pt, align=center,
    inner sep=8pt, text width=70mm,
    drop shadow={shadow xshift=0.4pt, shadow yshift=-0.7pt, opacity=0.22, fill=black!60}},
  harnessbox/.style={rounded corners=7pt, draw=black!30, line width=1pt, align=center,
    inner sep=8pt, text width=48mm,
    drop shadow={shadow xshift=0.4pt, shadow yshift=-0.7pt, opacity=0.22, fill=black!60}},
  flow/.style={->, black!55, line width=1.4pt},
  flowlab/.style={font=\sffamily\scriptsize\itshape, text=black!55},
]

% ================= middle column: the four stages =================
\node[stage, fill=envtint] (mine)
  {{\cardhd\color{minehd} 1.~Weakness Mining}\\[4pt]
   \begin{tikzpicture}[glyphrow]
     \pic at (0,0) {taskcard};
     \node at (0,-0.95) {run $\mathrm{H}_t$ on\\ held-in tasks};
     \draw[->, black!55] (0.6,0) -- (1.0,0);
     \pic at (1.7,0) {calltree};
     \node at (1.7,-0.95) {trace-based\\ diagnosis};
     \draw[->, black!55] (2.45,0) -- (2.85,0);
     \pic at (3.65,0) {clusterpair};
     \node at (3.65,-0.95) {cluster\\ by $\phi$};
   \end{tikzpicture}};

\node[stage, fill=editfill, below=12mm of mine] (prop)
  {{\cardhd\color{editgreen} 2.~Independent Candidate Proposal}\\[4pt]
   \begin{tikzpicture}[glyphrow]
     % stacked proposer inputs, each with its own arrow into the proposer card
     \node[anchor=east, align=right, font=\sffamily\scriptsize] (inp1) at (0.35,0.52) {operation evidence};
     \node[anchor=east, align=right, font=\sffamily\scriptsize] (inp2) at (0.35,0.1) {surface\\ capabilities};
     \node[anchor=east, align=right, font=\sffamily\scriptsize] (inp3) at (0.35,-0.32) {edit history};
     \foreach \n in {inp1,inp2,inp3}{
       \draw[->, black!55] ([xshift=2pt]\n.east) -- ([xshift=17pt]\n.east);}
     \pic at (1.55,0.1) {proposercard};
     \node[font=\scriptsize, text=black!55] at (1.55,-0.27) {$\mathrm{H}_t$};
     \node at (1.55,-0.8) {proposer\\(fixed model)};
     \draw[->, black!55] (2.15,0.1) -- (2.55,0.1);
     % k = 4 separate candidates, one edited surface each (never merged)
     \minigrid{2.75}{0.05}{4}{$C_1$}
     \minigrid{3.5}{0.05}{8}{$C_2$}
     \minigrid{4.25}{0.05}{4}{$C_3$}
     \minigrid{5.0}{0.05}{2}{$C_4$}
     \node at (3.92,-0.95) {$\le k$ separate harnesses,\\ one edited surface each};
   \end{tikzpicture}};

\node[stage, fill=scrtint, below=12mm of prop] (scr)
  {{\cardhd\color{scrhd} 3.~Screening Tournament}\\[4pt]
   \begin{tikzpicture}[glyphrow]
     \pic at (0,0.05) {subsetcards};
     \node at (0.05,-0.95) {$\mathrm{H}_t$, $C_1..C_k$ on small\\ subset $V_s$, $v_s$ attempt};
     \draw[->, black!55] (0.85,0.05) -- (1.25,0.05);
     % candidates ranked by metric change vs. the incumbent on V_s
     \draw[black!45, line width=0.5pt] (2.35,-0.48) -- (2.35,0.56);
     \node[rankbar, minimum width=1.05cm, fill=editgreen] at (2.35,0.36) {};
     \node[rankbar, minimum width=0.7cm, fill=black!35] at (2.35,0.12) {};
     \node[rankbar, minimum width=0.35cm, fill=black!35] at (2.35,-0.12) {};
     \node[rankbar, minimum width=0.5cm, fill=rejred!70, anchor=east] at (2.35,-0.36) {};
     \foreach \y/\c in {0.36/C_3,0.12/C_1,-0.12/C_4,-0.36/C_2}{
       \node[font=\sffamily\scriptsize, anchor=east] at (1.8,\y) {$\c$};}
     \node at (2.6,-0.95) {rank by $\Delta$\\ (ranking only)};
     \draw[->, black!55] (3.6,0.05) -- (4.0,0.05);
     \node[circle, fill=editgreen, text=white, inner sep=1.5pt, font=\sffamily\small]
       (win) at (4.55,0.05) {$C^\star$};
     \node at (4.55,-0.95) {top-ranked\\ advances};
   \end{tikzpicture}};

\node[stage, fill=subtint, below=12mm of scr] (val)
  {{\cardhd\color{valhd} 4.~Full Validation and Promotion}\\[4pt]
   \begin{tikzpicture}[glyphrow]
     \pic at (0,0) {splitcheck};
     \node at (0,-1.05) {$C^\star$ vs.\ $\mathrm{H}_t$ on\\ all of $V$, $v$ att.};
     \draw[->, black!55] (0.65,0) -- (1.35,0);
     % per-length F1 deltas against their veto lines -tau_l
     \draw[black!45, line width=0.5pt] (1.75,0.05) -- (3.45,0.05);
     \node[dbar, minimum height=0.32cm, fill=editgreen, anchor=south] at (2.05,0.05) {};
     \node[dbar, minimum height=0.14cm, fill=rejred!70, anchor=north] at (2.6,0.05) {};
     \node[dbar, minimum height=0.48cm, fill=editgreen, anchor=south] at (3.15,0.05) {};
     \foreach \x/\t in {2.05/0.38,2.6/0.42,3.15/0.55}{
       \draw[rejred, densely dashed, line width=0.6pt] (\x-0.2,0.05-\t) -- (\x+0.2,0.05-\t);}
     \foreach \x/\l in {2.05/8K,2.6/16K,3.15/32K}{
       \node[font=\sffamily\scriptsize, text=black!55] at (\x,0.72) {\l};}
     \node at (2.6,-1.05) {$\Delta_\ell\ge-\tau_\ell$ at every length,\\ macro $\ge\tau_{\mathrm{imp}}$, cost band};
     \draw[->, black!55] (3.85,0) -- (4.45,0);
     \pic at (4.85,0) {acceptreject};
     \node at (4.95,-1.05) {promote /\\ reject};
   \end{tikzpicture}};

% ================= side columns: current and updated harness =================
\coordinate (sidemid) at ($(prop.south)!0.5!(scr.north)$);
\node[harnessbox, fill=fixfill, anchor=east] (cur) at ([xshift=-10mm]sidemid -| prop.west)
  {{\cardhd Current RLM Harness $\mathrm{H}_t$~\tikz[baseline=-2pt, scale=0.8,
     sharp corners]{% sharp corners: the box node's rounded corners=7pt leaks
     % into this inline picture and would distort the small glyph paths
     \fill[inkcol] (0,0.22) circle (0.11);
     \fill[inkcol] (-0.2,-0.18) arc (180:0:0.2) -- cycle;}}\\[4pt]
   \begin{tikzpicture}[glyphrow]
     % all placement via node anchors: raw-coordinate paths drift relative to
     % nodes in this document, but paths anchored on node coordinates land true
     \node[rounded corners=2pt, draw=fixgrey, dashed, line width=0.7pt, fill=white,
       minimum width=1.06cm, minimum height=1.0cm, align=center] (mc) at (0,0) {fixed\\ model};
     \path (mc.north east) pic[scale=1.3] {lockbadge};
     \foreach \j [count=\r from 0] in {2,1,0}{
       \foreach \i in {0,1,2}{
         \pgfmathtruncatemacro{\slab}{\r*3+\i+1}
         \ifnum\slab=2 \tikzset{c/.style=celledited}\else
         \ifnum\slab=6 \tikzset{c/.style=celledited}\else
           \tikzset{c/.style=cellplain}\fi\fi
         \node[c] (curcell\slab) at ({1.62+\i*0.46},{\j*0.46-0.46}) {S\slab};}}
     % + drawn as path strokes anchored on node coordinates (a text node here
     % drifts off the true midpoint)
     \coordinate (plusmid) at ($(mc.east)!0.5!(curcell4.west)$);
     \draw[black!55, line width=1.1pt, line cap=round]
       ($(plusmid)+(-0.12,0)$) -- ($(plusmid)+(0.12,0)$);
     \draw[black!55, line width=1.1pt, line cap=round]
       ($(plusmid)+(0,-0.12)$) -- ($(plusmid)+(0,0.12)$);
     \node[cellplain] (curcell10) at ({1.62+1*0.46},{-1*0.46-0.46}) {S10};
     \node[anchor=north] (eslab) at ([yshift=-4.5pt]curcell10.south) {Editable Surfaces};
     \node[anchor=north, font=\sffamily\scriptsize, text=black!55] at ([yshift=-1pt]eslab.south) {green: prev. edited};
   \end{tikzpicture}};

\node[harnessbox, fill=editfill, anchor=west] (new) at ([xshift=10mm]sidemid -| prop.east)
  {{\cardhd Updated RLM Harness $\mathrm{H}_{t+1}$~\tikz[baseline=-2pt, scale=0.8,
     sharp corners]{% sharp corners: see the current-harness title note
     \fill[inkcol] (0,0.22) circle (0.11);
     \fill[inkcol] (-0.2,-0.18) arc (180:0:0.2) -- cycle;}}\\[4pt]
   \begin{tikzpicture}[glyphrow]
     \node[rounded corners=2pt, draw=fixgrey, dashed, line width=0.7pt, fill=white,
       minimum width=1.06cm, minimum height=1.0cm, align=center] (mcn) at (0,0) {fixed\\ model};
     \path (mcn.north east) pic[scale=1.3] {lockbadge};
     \foreach \j [count=\r from 0] in {2,1,0}{
       \foreach \i in {0,1,2}{
         \pgfmathtruncatemacro{\slab}{\r*3+\i+1}
         \ifnum\slab=2 \tikzset{c/.style=celledited}\else
         \ifnum\slab=4 \tikzset{c/.style=cellnew}\else
         \ifnum\slab=6 \tikzset{c/.style=celledited}\else
           \tikzset{c/.style=cellplain}\fi\fi\fi
         \node[c] (newcell\slab) at ({1.62+\i*0.46},{\j*0.46-0.46}) {S\slab};}}
     \coordinate (plusmidn) at ($(mcn.east)!0.5!(newcell4.west)$);
     \draw[black!55, line width=1.1pt, line cap=round]
       ($(plusmidn)+(-0.12,0)$) -- ($(plusmidn)+(0.12,0)$);
     \draw[black!55, line width=1.1pt, line cap=round]
       ($(plusmidn)+(0,-0.12)$) -- ($(plusmidn)+(0,0.12)$);
     \node[cellplain] (newcell10) at ({1.62+1*0.46},{-1*0.46-0.46}) {S10};
   \end{tikzpicture}\\[4pt]
   {\cardsb $\mathrm{H}_t$ plus the winner's single edit (\sid{S4})}};

% ================= loop arrows =================
% current harness up into weakness mining
\draw[flow] (cur.north) |- node[flowlab, pos=0.72, above]{run on tasks} (mine.west);
% mining -> proposal -> screening -> validation (main path, left of centre)
\draw[flow] ([xshift=-14mm]mine.south) -- node[flowlab, left=1mm]{failure patterns} ([xshift=-14mm]prop.north);
\draw[flow] ([xshift=-14mm]prop.south) -- node[flowlab, left=1mm, align=right]{preflight-admitted\\ candidates} ([xshift=-14mm]scr.north);
\draw[flow] ([xshift=-14mm]scr.south) -- node[flowlab, left=1mm]{winner $C^\star$ only} ([xshift=-14mm]val.north);
% fallback 1: every screened candidate rejected -> re-attribute the same mining
% runs once and propose a fresh pool
\draw[flow, dashed, scrhd] ([xshift=14mm]scr.north) -- node[flowlab, text=scrhd, right=1mm, align=left]{all screen-rejected:\\ re-attribute once} ([xshift=14mm]prop.south);
% fallback 2: rejected winner with a conditional-F1 gain -> one repair generation
\draw[flow, dashed, valhd] ([xshift=-10mm]val.south) -- ++(0,-6mm) -| node[flowlab, text=valhd, pos=0.25, below, align=center]{rejected but conditional F1 gain: mine $C^\star$, repair once, re-validate vs.\ $\mathrm{H}_t$} ([xshift=10mm]val.south);
% rejected: back to the same harness
\draw[flow, rejred] (val.west) -| node[flowlab, text=rejred, pos=0.25, below, align=center]{reject --- keep $\mathrm{H}_t$,\\ retire other candidates} (cur.south);
% promotion adopts the measured single-edit winner
\draw[flow, editgreen] (val.east) -| node[flowlab, text=editgreen, pos=0.25, below]{promote $C^\star$} (new.south);
% updated harness starts the next round
\draw[flow, dashed] (new.north) -- ++(0,12.5mm) node[flowlab, above=0.5mm, align=center]{next round\\ ($\mathrm{H}_{t+1}$ becomes current)};

\end{tikzpicture}}
\caption{One \sh{} round. Final-answer outcomes and selected execution operations
support model-inferred diagnoses. The proposer returns up to $k$ candidates, each
a separate harness with one edited surface; candidates are never merged. A
screening tournament runs the incumbent and every candidate on a small subset
$V_s$ of the held-out tasks and only ranks them. The top-ranked $C^\star$ alone is
compared with the incumbent on all held-out tasks; it must clear a per-length F1
regression veto $\tau_\ell$, a macro-F1 improvement threshold, and the
wrong-format failure veto. Cost is recorded for the later audit but is not an
enforced gate in the reported run.
Dashed paths are bounded fallbacks: one re-attribution of the same mining runs
when every screened candidate is rejected, and one repair generation when a
rejected winner improves F1 on runs with a measurable answer. Validation traces
are never mined. Side grids illustrate previously edited surfaces and one promoted
edit. Weights and verifier stay fixed.}
\label{fig:sh-loop}
\end{figure*}
